% !TEX program = xelatex
\documentclass[11pt]{article}

\usepackage[a4paper,margin=1in]{geometry}
\usepackage{amsmath,amssymb}
\usepackage{booktabs}
\usepackage{enumitem}
\usepackage{hyperref}
\usepackage{xcolor}
\usepackage{kotex}
\usepackage{listings}

\lstdefinestyle{codestyle}{
    basicstyle=\ttfamily\small,
    breaklines=true,
    frame=single,
    columns=fullflexible,
    backgroundcolor=\color{gray!5}
}
\lstset{style=codestyle}

\title{Phase 1 DeepXDE Code Guide\\
\large Robot Joint Motor Housing Inverse PINN}
\author{안형준}
\date{}

\begin{document}
\maketitle

\section{문서 목적}

이 문서는 \texttt{main\_deepxde.py}의 코드 흐름을 설명하기 위한 해설 문서다.

\begin{center}
\begin{tabular}{ll}
\toprule
역할 & 파일 \\
\midrule
비교용 PyTorch baseline & \texttt{main.py} \\
DeepXDE 버전 PINN 코드 & \texttt{main\_deepxde.py} \\
기존 코드 설명 문서 & \texttt{code\_guide/PHASE1\_CODE\_GUIDE.tex} \\
\bottomrule
\end{tabular}
\end{center}

\texttt{main\_deepxde.py}는 \texttt{main.py}를 지우거나 대체하지 않고, 같은 inverse PINN 문제를 DeepXDE 형식으로 따로 구현한 파일이다. 현재 버전은 \texttt{main.py}에서 값을 import하지 않고, 필요한 물리 상수와 helper 함수를 파일 안에 직접 정의한다.

\section{한 줄 요약}

DeepXDE 버전도 푸는 물리 문제는 같다.

\[
\text{input} = (X,\tau),
\qquad
\text{output} = \theta(X,\tau)
\]

학습 대상 inverse parameter는 다음 두 개다.

\[
\beta \rightarrow h,
\qquad
S \rightarrow q_0
\]

차이는 loss를 직접 PyTorch 함수로 하나씩 계산하지 않고, 다음 DeepXDE 객체에 맡긴다는 점이다.

\begin{center}
\begin{tabular}{ll}
\toprule
DeepXDE 객체 & 역할 \\
\midrule
\texttt{TimePDE} & PDE residual과 학습점 관리 \\
\texttt{IC} & 초기조건 관리 \\
\texttt{OperatorBC} & 미분 경계조건 관리 \\
\texttt{PointSetBC} & 센서 데이터 조건 관리 \\
\bottomrule
\end{tabular}
\end{center}

\section{실행 방법}

현재 코드에서는 DeepXDE backend로 PyTorch를 사용한다.

\begin{lstlisting}[language=bash]
python3 -m pip install deepxde
python3 main_deepxde.py --iterations 5000
\end{lstlisting}

짧게 구조만 확인하려면 다음처럼 실행할 수 있다.

\begin{lstlisting}[language=bash]
python3 main_deepxde.py --iterations 10 --display-every 1
\end{lstlisting}

\section{PyTorch baseline과 DeepXDE 코드의 대응}

아래 표는 비교를 위한 대응 관계다. 실제 \texttt{main\_deepxde.py}는 \texttt{main.py}를 import하지 않고 독립적으로 실행된다.

\begin{center}
\begin{tabular}{lll}
\toprule
역할 & \texttt{main.py} & \texttt{main\_deepxde.py} \\
\midrule
신경망 & \texttt{MotorHousingPINN} & \texttt{dde.nn.FNN} \\
PDE residual & \texttt{pde\_residual()} & \texttt{pde()} \\
PDE 점 샘플링 & \texttt{sample\_collocation\_points()} & \texttt{dde.data.TimePDE} \\
초기조건 & \texttt{initial\_condition\_loss()} & \texttt{dde.icbc.IC} \\
단열 경계조건 & \texttt{boundary\_condition\_loss()} & \texttt{dde.icbc.OperatorBC} \\
센서 데이터 조건 & \texttt{sensor\_data\_loss()} & \texttt{dde.icbc.PointSetBC} \\
학습 파라미터 & \texttt{raw\_beta}, \texttt{raw\_s} & \texttt{dde.Variable} \\
학습 실행 & 아직 smoke test 중심 & \texttt{model.compile()}, \texttt{model.train()} \\
\bottomrule
\end{tabular}
\end{center}

\section{물리 상수와 함수의 독립 정의}

\texttt{main\_deepxde.py}는 \texttt{main.py}에서 값을 import하지 않는다. DeepXDE 코드 파일 하나만 봐도 실행 조건을 이해할 수 있도록 형상값, 물성치, 무차원화 기준값, true parameter, helper 함수를 파일 안에 직접 정의한다.

코드 상단에는 \texttt{math}도 import한다. 이는 단면적 \(A\)와 둘레 \(P\)를 계산할 때 \(\pi\)가 필요하기 때문이다.

\begin{lstlisting}[language=Python]
import math
\end{lstlisting}

\begin{lstlisting}[language=Python]
L = 0.15
R_OUTER = 0.04
R_INNER = 0.025
X0 = 0.30
SIGMA_X = 0.10
SENSOR_X = (0.25, 0.55, 0.85)

RHO = 2700.0
CP = 900.0
K = 167.0
T_INF = 25.0

DELTA_T_REF = 50.0
I_REF = 5.0
T_FINAL = 1200.0

H_TRUE = 40.0
Q0_TRUE = 8.0e4
\end{lstlisting}

입력값에서 자동 계산되는 값들도 같은 파일 안에서 계산한다.

\begin{lstlisting}[language=Python]
ALPHA = K / (RHO * CP)
AREA = math.pi * (R_OUTER**2 - R_INNER**2)
PERIMETER = 2.0 * math.pi * R_OUTER
TAU_FINAL = ALPHA * T_FINAL / L**2
\end{lstlisting}

\(\beta_{\mathrm{true}}\)와 \(S_{\mathrm{true}}\)는 변환 함수가 정의된 뒤 같은 파일 안에서 계산한다.

\begin{lstlisting}[language=Python]
BETA_TRUE = beta_from_physical_h(H_TRUE)
S_TRUE = s_from_physical_q0(Q0_TRUE)
\end{lstlisting}

다음 helper 함수들도 \texttt{main\_deepxde.py} 안에 직접 들어 있다.

\begin{center}
\begin{tabular}{ll}
\toprule
함수 & 역할 \\
\midrule
\texttt{tau\_from\_time()} & 실제 시간 \(t\)를 무차원 시간 \(\tau\)로 변환 \\
\texttt{time\_from\_tau()} & 무차원 시간 \(\tau\)를 실제 시간 \(t\)로 변환 \\
\texttt{beta\_from\_physical\_h()} & \(h\)를 \(\beta\)로 변환 \\
\texttt{s\_from\_physical\_q0()} & \(q_0\)를 \(S\)로 변환 \\
\texttt{gaussian\_heat\_source()} & \(g(X)\) 계산 \\
\texttt{current\_profile()} & \(i(\tau)\) 계산 \\
\texttt{physical\_h\_from\_beta()} & \(\beta\)를 \(h\)로 변환 \\
\texttt{physical\_q0\_from\_s()} & \(S\)를 \(q_0\)로 변환 \\
\texttt{synthetic\_theta\_no\_conduction()} & local ODE 기반 synthetic data 생성 \\
\bottomrule
\end{tabular}
\end{center}

이 방식은 \texttt{main.py}와 코드 중복이 생기지만, DeepXDE 버전이 독립적인 실행 파일이 된다는 장점이 있다.

\subsection{파일 내부 구역}

\texttt{main\_deepxde.py}의 큰 구역은 다음 순서로 되어 있다.

\begin{center}
\begin{tabular}{ll}
\toprule
구역 & 역할 \\
\midrule
1 & 사용자가 정하는 입력값 \\
2 & 입력값에서 자동 계산되는 값 \\
3 & 차원/무차원 변환 함수 \\
4 & heat source, current profile, synthetic data 함수 \\
5 & inverse parameter 양수 제약 \\
6 & DeepXDE용 sensor data 생성 \\
7 & command line option \\
8 & DeepXDE model 구성 \\
9 & 출력 helper \\
10 & 실행부 \\
\bottomrule
\end{tabular}
\end{center}

\section{DeepXDE backend와 device 설정}

파일 상단에서 DeepXDE backend를 PyTorch로 지정한다.

\begin{lstlisting}[language=Python]
os.environ.setdefault("DDE_BACKEND", "pytorch")
os.environ.setdefault("MPLCONFIGDIR", "/private/tmp/matplotlib-cache")
\end{lstlisting}

그래서 \texttt{pde()} 함수 안에서는 \texttt{torch.Tensor}와 PyTorch 연산을 그대로 사용할 수 있다.

추가로 현재 코드는 기본 실행 device를 CPU로 고정한다.
DeepXDE 1.15의 PyTorch backend는 Mac에서 MPS가 가능하면 import 시점에 \texttt{torch.set\_default\_device("mps")}를 호출한다.
M1 Mac과 VSCode Jupyter 조합에서는 이 MPS 자동 선택이 \texttt{metal gpu stream} crash를 만들 수 있으므로, 기본값은 CPU 실행으로 둔다.

\begin{lstlisting}[language=Python]
FORCE_CPU = os.environ.get("PINN_FORCE_CPU", "1") != "0"

if FORCE_CPU:
    torch.set_default_device("cpu")
    torch.backends.mps.is_available = _disable_mps_for_deepxde
\end{lstlisting}

MPS를 다시 실험하고 싶으면 실행 전에 \texttt{PINN\_FORCE\_CPU=0}을 지정하면 된다.

\section{양수 파라미터 처리}

물리적으로 \(\beta\), \(S\), \(h\), \(q_0\)는 음수가 되면 이상하다. 따라서 \texttt{main\_deepxde.py}에서도 raw parameter를 학습하고, softplus를 통과시켜 양수로 만든다.

\begin{lstlisting}[language=Python]
def positive_parameter(raw_value):
    return F.softplus(raw_value)
\end{lstlisting}

\[
\beta = \mathrm{softplus}(\texttt{raw\_beta})
\]

\[
S = \mathrm{softplus}(\texttt{raw\_s})
\]

\texttt{raw\_beta}와 \texttt{raw\_s}는 DeepXDE에서 다음처럼 정의한다.

\begin{lstlisting}[language=Python]
raw_beta = dde.Variable(0.0)
raw_s = dde.Variable(0.0)
\end{lstlisting}

초기값이 0이면 softplus 때문에 실제 초기 추정값은 다음과 같다.

\[
\mathrm{softplus}(0) \approx 0.693
\]

\section{Synthetic sensor data}

센서 데이터는 \texttt{main\_deepxde.py} 안에 직접 정의된 \texttt{solve\_synthetic\_theta\_finite\_difference()}로 만든다.
즉 전도항을 제외한 local ODE가 아니라, 전체 PDE를 finite difference method로 푼 synthetic data를 사용한다.
다만 DeepXDE의 \texttt{PointSetBC}는 NumPy 배열을 받기 때문에 \texttt{torch.Tensor}를 NumPy로 변환한다.

finite difference solver는 MPS crash를 피하기 위해 내부 tensor를 명시적으로 CPU에 생성한다.
마지막 변환에서도 \texttt{.detach().cpu()}를 호출해서 계산 그래프에서 분리하고 CPU memory의 NumPy 배열로 바꾼다.

\begin{lstlisting}[language=Python]
def make_synthetic_sensor_data_numpy(n_time=41):
    xtau, theta_sensor = solve_synthetic_theta_finite_difference(n_time)
    return (
        xtau.detach().cpu().numpy().astype("float32"),
        theta_sensor.detach().cpu().numpy().astype("float32"),
    )
\end{lstlisting}

기본 센서 데이터 개수는 다음과 같다.

\[
3~\text{sensors} \times 41~\text{time points}
=
123~\text{points}
\]

주의할 점은 현재 synthetic data가 full PDE 해가 아니라는 것이다. 아직 전도항 \(\theta_{XX}\)가 없는 local ODE 기반 데이터다.

\[
\frac{d\theta}{d\tau}
=
-\beta_{\mathrm{true}}\theta
+
S_{\mathrm{true}} i(\tau)^2 g(X)
\]

\section{DeepXDE PDE residual}

제안서 기준 무차원 PDE는 다음과 같다.

\[
\frac{\partial \theta}{\partial \tau}
=
\frac{\partial^2 \theta}{\partial X^2}
-
\beta\theta
+
S i(\tau)^2 g(X)
\]

residual 형태는 다음과 같다.

\[
f_{\mathrm{PDE}}
=
\theta_\tau
-
\theta_{XX}
+
\beta\theta
-
S i(\tau)^2 g(X)
\]

\texttt{main\_deepxde.py}에서는 이 식이 \texttt{pde()} 함수에 들어간다.

\begin{lstlisting}[language=Python]
def pde(xtau, theta):
    theta_tau = dde.grad.jacobian(theta, xtau, i=0, j=1)
    theta_xx = dde.grad.hessian(theta, xtau, component=0, i=0, j=0)

    x = xtau[:, 0:1]
    tau = xtau[:, 1:2]

    beta = positive_parameter(raw_beta)
    s_value = positive_parameter(raw_s)
    i_tau = current_profile(tau)
    g_x = gaussian_heat_source(x)

    return theta_tau - theta_xx + beta * theta - s_value * (i_tau**2) * g_x
\end{lstlisting}

여기서 \texttt{xtau[:, 0:1]}은 \(X\), \texttt{xtau[:, 1:2]}는 \(\tau\)다.

\section{해석 영역}

DeepXDE에서는 공간 영역과 시간 영역을 따로 만든 뒤 합친다.

\begin{lstlisting}[language=Python]
geom = dde.geometry.Interval(0.0, 1.0)
time_domain = dde.geometry.TimeDomain(0.0, TAU_FINAL)
geomtime = dde.geometry.GeometryXTime(geom, time_domain)
\end{lstlisting}

즉 해석 영역은 다음과 같다.

\[
0 \le X \le 1,
\qquad
0 \le \tau \le \tau_f
\]

\section{초기조건}

초기조건은 다음과 같다.

\[
\theta(X,0)=0
\]

DeepXDE에서는 다음처럼 정의한다.

\begin{lstlisting}[language=Python]
ic = dde.icbc.IC(geomtime, zero_numpy, on_initial)
\end{lstlisting}

\texttt{zero\_numpy()}는 초기 온도 \(\theta=0\)을 반환하는 함수다.

\section{단열 경계조건}

제안서의 양끝 경계조건은 단열이다.

\[
\frac{\partial \theta}{\partial X}(0,\tau)=0
\]

\[
\frac{\partial \theta}{\partial X}(1,\tau)=0
\]

DeepXDE에서는 미분 조건을 \texttt{OperatorBC}로 넣는다.

\begin{lstlisting}[language=Python]
def adiabatic_bc(inputs, outputs, _):
    return dde.grad.jacobian(outputs, inputs, i=0, j=0)

left_bc = dde.icbc.OperatorBC(geomtime, adiabatic_bc, on_left_boundary)
right_bc = dde.icbc.OperatorBC(geomtime, adiabatic_bc, on_right_boundary)
\end{lstlisting}

여기서 \texttt{j=0}은 입력 \((X,\tau)\) 중 첫 번째 변수인 \(X\)에 대한 미분이라는 뜻이다.

\section{센서 데이터 조건}

센서 데이터는 DeepXDE의 \texttt{PointSetBC}로 넣는다.

\begin{lstlisting}[language=Python]
sensor_bc = dde.icbc.PointSetBC(sensor_xtau, sensor_theta, component=0)
\end{lstlisting}

이 조건은 센서 위치와 시간에서 신경망 예측값이 synthetic sensor data와 가까워지도록 만든다.

\[
L_{\mathrm{data}}
=
\frac{1}{N_s}
\sum_i
\left[
\theta_\phi(X_i,\tau_i)-\theta_{\mathrm{sensor},i}
\right]^2
\]

\section{TimePDE 데이터 객체}

DeepXDE의 \texttt{TimePDE}에 PDE, 초기조건, 경계조건, 센서 데이터 조건을 한 번에 넘긴다.

\begin{lstlisting}[language=Python]
data = dde.data.TimePDE(
    geomtime,
    pde,
    [ic, left_bc, right_bc, sensor_bc],
    num_domain=args.num_domain,
    num_boundary=args.num_boundary,
    num_initial=args.num_initial,
    num_test=args.num_test,
    train_distribution="Hammersley",
)
\end{lstlisting}

각 인자의 의미는 다음과 같다.

\begin{center}
\begin{tabular}{ll}
\toprule
인자 & 의미 \\
\midrule
\texttt{num\_domain} & PDE residual을 검사할 내부 점 개수 \\
\texttt{num\_boundary} & 경계조건 점 개수 \\
\texttt{num\_initial} & 초기조건 점 개수 \\
\texttt{num\_test} & 테스트용 점 개수 \\
\texttt{train\_distribution} & 학습점 샘플링 방식 \\
\bottomrule
\end{tabular}
\end{center}

\section{신경망 구조}

DeepXDE에서는 신경망을 다음처럼 만든다.

\begin{lstlisting}[language=Python]
layer_sizes = [2] + [args.width] * args.depth + [1]
net = dde.nn.FNN(layer_sizes, "tanh", "Glorot normal")
\end{lstlisting}

기본값은 다음과 같다.

\[
[2, 32, 32, 32, 1]
\]

즉 입력은 \(X,\tau\) 두 개이고, 출력은 \(\theta\) 하나다.

\section{학습 실행}

\texttt{main()}에서는 먼저 Adam optimizer로 학습한다.

\begin{lstlisting}[language=Python]
model.compile(
    "adam",
    lr=args.lr,
    external_trainable_variables=[raw_beta, raw_s],
)
model.train(iterations=args.iterations, display_every=args.display_every)
\end{lstlisting}

중요한 부분은 \texttt{external\_trainable\_variables}다. DeepXDE의 신경망 weight뿐 아니라 \(\beta\), \(S\)에 해당하는 \texttt{raw\_beta}, \texttt{raw\_s}도 같이 학습시키라는 뜻이다.

옵션으로 L-BFGS 후처리도 가능하다.

\begin{lstlisting}[language=Python]
python3 main_deepxde.py --iterations 5000 --lbfgs
\end{lstlisting}

\section{학습 결과 출력}

학습이 끝나면 \(\beta\), \(S\)를 실제 물리 파라미터로 바꿔서 출력한다.

\begin{lstlisting}[language=Python]
beta = positive_parameter(raw_beta).detach()
s_value = positive_parameter(raw_s).detach()
h_value = physical_h_from_beta(beta)
q0_value = physical_q0_from_s(s_value)
\end{lstlisting}

변환식은 \texttt{main\_deepxde.py} 안에 직접 정의된 다음 두 함수가 사용하는 식이다.

\begin{lstlisting}[language=Python]
physical_h_from_beta()
physical_q0_from_s()
\end{lstlisting}

\[
h =
\frac{\beta kA}{PL^2}
\]

\[
q_0 =
\frac{S k\Delta T_{\mathrm{ref}}}{I_{\mathrm{ref}}^2L^2}
\]

\section{명령행 인자}

\begin{center}
\begin{tabular}{lll}
\toprule
인자 & 기본값 & 의미 \\
\midrule
\texttt{--iterations} & 5000 & Adam 학습 반복 횟수 \\
\texttt{--lr} & \(1.0\times10^{-3}\) & learning rate \\
\texttt{--depth} & 3 & hidden layer 개수 \\
\texttt{--width} & 32 & hidden layer neuron 개수 \\
\texttt{--num-domain} & 2560 & PDE 내부 점 개수 \\
\texttt{--num-boundary} & 256 & 경계 점 개수 \\
\texttt{--num-initial} & 128 & 초기조건 점 개수 \\
\texttt{--num-test} & 1000 & 테스트 점 개수 \\
\texttt{--sensor-time-points} & 41 & 센서별 시간점 개수 \\
\texttt{--display-every} & 500 & loss 출력 간격 \\
\texttt{--seed} & 0 & random seed \\
\texttt{--lbfgs} & off & Adam 이후 L-BFGS 실행 여부 \\
\bottomrule
\end{tabular}
\end{center}

\section{현재 DeepXDE 코드의 한계}

현재 \texttt{main\_deepxde.py}가 아직 하지 않는 것은 다음과 같다.

\begin{enumerate}[leftmargin=2em]
\item full PDE finite difference solver 기반 synthetic data 생성
\item loss weight tuning
\item temperature contour plot
\item sensor fitting plot
\item parameter convergence plot
\item 실제 실험 데이터 입력
\end{enumerate}

가장 중요한 한계는 sensor data가 아직 전도항을 포함한 full PDE data가 아니라 local ODE 기반이라는 점이다.

\[
\theta_{XX} \quad \text{is not included in current synthetic data}
\]

따라서 현재 DeepXDE 코드는 최종 검증용이라기보다, DeepXDE로 inverse PINN 구조를 구현하고 학습 흐름을 확인하는 버전으로 보는 것이 좋다.

\section{다음 구현 순서}

추천 순서는 다음과 같다.

\begin{enumerate}[leftmargin=2em]
\item DeepXDE 설치 후 짧은 iteration으로 실행 확인
\item \(\beta\), \(S\)가 학습 중 변하는지 확인
\item loss weight를 조정해서 data loss와 PDE loss 균형 확인
\item finite difference solver로 full PDE synthetic data 생성
\item DeepXDE 결과와 PyTorch baseline 결과 비교
\item learned \(h\), \(q_0\)를 true value와 비교하는 plot 추가
\end{enumerate}

\end{document}
